% Generated by scripts/build_jss_manuscript.py from Phase 29.

\documentclass[article,nojss]{jss}

\usepackage{lmodern}

\setlength{\emergencystretch}{2em}

\author{Soyeb Mohammad}

\Plainauthor{Soyeb Mohammad}

\Address{Soyeb Mohammad\\Author affiliation and correspondence pending confirmation.}

\title{When Sampling Frames Share a Blind Spot: A Reproducible Benchmark for Population-Mean Certification}

\Plaintitle{When sampling frames share a blind spot: a reproducible benchmark for population-mean certification}

\Shorttitle{Certification under Shared Undercoverage}

\Abstract{Agreement between sampling channels can leave both channels blind to the same target units. We present GENEVRA, a reproducible benchmark that separates target-mean coverage, interval precision, threshold decisions, and anchor certification under controlled information and costs. In a frozen shared-blindness case, 0 of 512 intervals covered the target and all 180 issued certificates were false, although no threshold label was wrong. In a separate matched-cost comparison at 32,768 charged units, GENEVRA's width-proxy allocation gave 1 of 256 covered-and-precise intervals, versus 104 for paid-pilot cost-adjusted Neyman and 229 for equal spend. The paired proxy-minus-Neyman difference was -40.23 percentage points (95\% Hoeffding interval -57.21 to -23.26). In a separate large-omission case, applying a valid supplied cap to the same equal-spend acquisitions changed coverage from zero to all runs, with no covered-and-precise intervals. An independent implementation of the numerical calculations, executed internally, recovered the central results. The contribution is an auditable benchmark of these failure modes. Its evidence concerns synthetic populations under stated laws, restrictions and costs.}

\Keywords{undercoverage, partial identification, stratified allocation, statistical certification, reproducible benchmarking}

\Plainkeywords{undercoverage, partial identification, stratified allocation, statistical certification, reproducible benchmarking}

\begin{document}

\section{Introduction}

Sampling precision and target identification answer different questions.
A large calibration sample can estimate the distribution on a frame
accurately while leaving the population outside that frame unobserved.
Two separately sampled channels may agree because they share omissions.
An audit that compares only those channels can then admit a policy that
reports a precise but wrong target interval or certifies an anchor whose
error exceeds the intended tolerance.

The resulting failures can look similar in aggregate. A threshold
decision can remain correct when the interval misses the target on the
same side of the threshold. A policy can achieve coverage by returning a
broad interval, or reduce false certificates by withholding them. A
cost-aware allocator can correctly optimize a surrogate loss and still
produce fewer useful intervals. The benchmark keeps these outcomes
separate so that method comparisons can distinguish sources of failure.

\subsection{Research questions and contribution}

We ask when shared omission produces false certificates despite
agreement, when defensible external restrictions recover target
coverage, and how allocation changes useful precision at matched charged
cost. A fourth question concerns auditability. Can separate code recover
the decisions and comparisons from underlying records? The contribution
is the linked suite of frozen stress cases, access and tariff controls,
raw-count verification and negative comparisons. Missing-data bounds,
confidence intervals and stratified allocation supply the established
statistical methods.

\section{Related work and positioning}

Manski \citep{manski2005} and Imbens and Manski \citep{imbens2004}
distinguish uncertainty about an identified set from sampling
uncertainty about the parameter and its bounds. Tudball and colleagues
\citep{tudball2023} show how auxiliary constraints can inform inference
under selection bias. Their application estimates constraints in a
different model. GENEVRA supplies restrictions as controlled inputs,
including restrictions that are deliberately wrong, and measures their
effect on decisions.

Multiple-frame survey methods use probability designs and overlap
information to extend access \citep{lohr2006}. A spatial probability
survey of Norwegian recreational fisheries exposed coverage gaps in
alternative frames \citep{volstad2020}. This is a field precedent for
independently checking coverage. GENEVRA contributes synthetic known
truth and matched costs rather than new field evidence. Independent
draws from incomplete frames do not supply information about units
excluded from both.

Selective classification and held-out calibration distinguish the
decisions issued from their errors \citep{geifman2017,angelopoulos2025}.
Our object is a population mean relative to an anchor, with correctness
evaluated against declared synthetic truth. Neyman allocation and exact
cost-constrained allocation provide variance objectives
\citep{neyman1934,wright2019}. The width proxy tested here has a
different objective. Following the simulation-design distinctions in
Morris and colleagues \citep{morris2019}, we fix the generating laws,
estimands, methods and measures, then separate prespecified contrasts
from later sensitivity analyses. The closest prior work limits the
originality claim to the specific experimental and verification
artifact.

\section{Mathematical formulation}

\begin{table}[htbp]
\centering
\small
\begin{tabular}{lp{0.64\linewidth}}
\hline
Symbol & Meaning \\
\hline
T, \(F_A\), \(F_B\) & Target population and two sampling frames \\
\(w_h\), \(p_h\) & Reachable stratum weight and Bernoulli tail rate \\
r, m, z & Reachable mixture rate, omitted mass, hidden tail rate \\
U, V & Externally supplied upper bounds on omitted mass and hidden tail rate \\
\(c_h\), \(n_h\), B & Unit price, independent-unit count, charged acquisition budget \\
a, \ensuremath{\delta}, \ensuremath{\epsilon}, \ensuremath{\tau} & Anchor score, certificate tolerance, decision threshold, width limit \\
\(I=[L,H]\), \ensuremath{\theta} & Score interval and target score mean \\
\hline
\end{tabular}
\caption{Core notation. Phase 18 and Phase 25 share the score cutoffs but use distinct certificate constructions.}
\label{tab:notation}
\end{table}

\subsection{Population, frames and information}

Let T denote the target and \(F_A\) and \(F_B\) the subsets reachable by
two sampling channels. Shared undercoverage is \(T\setminus(F_A\cup
F_B)\). Asymmetric undercoverage excludes a domain in one channel that
the other can reach. Independent random draws conditional on each frame
say nothing about whether the exclusions overlap. In the frozen
shared-blind case the channels observe the source law despite a changed
evaluation law. This is a planted common blind spot. The scenario does
not estimate the prevalence of correlated omissions.

The reduced allocation model has two reachable strata with known
conditional weights \(w_h\) summing to one. Outcomes within a stratum
are independent Bernoulli tail indicators with rates \(p_h\). Set
\(r=\sum_h w_h p_h\) and \(q_T=(1-m)r+mz\). The target mean is
\(\theta=0.30+234.2q_T\). The conditional body and tail score means are
supplied constants. Neither method samples the missing population. The
evaluator knows \(q_T\) and latent rates for scoring, and operational
policies receive the same declared inputs. Only the labeled oracle
allocates using true stratum variances.

\subsection{Intervals and identification}

After an independent pilot, stratum h has \(k_h\) successes in \(n_h\)
independent inference draws. Its point rate is \(k_h/n_h\). For the
fixed-look likelihood bounds, retain rates p with log likelihood minus
its maximized value at least \(\log(\alpha/4)\). The factor four
accounts for two tails and two strata. Weight the lower and upper
endpoints separately to obtain the reachable band \([l,u]\). A zero
observed count gives an analytic upper endpoint, and an all-success
count gives the corresponding lower endpoint. The nominal interpretation
requires the stated sampling law, weights, and independence of pilot and
inference.

Given supplied bounds \(0\leq m\leq U\) and \(0\leq z\leq V\), the
target rate hull is \([(1-U)l,u+U\max(V-u,0)]\). This follows by taking
endpoints of the affine expression in m and z for each reachable
endpoint. Apply the affine score transformation to the hull. With an
exactly known r and \(V=1\) the width is \(234.2U\). Reachable sampling
cannot remove that identification component. Unknown or invalid
restrictions prevent the same interpretation for the full target.

\subsection{Coverage, precision and certificates}

A score interval \([L,H]\) covers when \(L\leq\theta\leq H\). Precision
means \(H-L\leq\tau\). Useful precision is their intersection. A
beneficial label requires \(L>\epsilon\). A nonbeneficial label requires
\(H\leq\epsilon\). Otherwise the policy abstains from threshold
classification. The frozen \ensuremath{\epsilon} and \ensuremath{\tau}
are both 0.60, in different roles. The anchor tolerance
\ensuremath{\delta} is 0.30 score units.

In Phase 25 a certificate is issued when the entire interval lies inside
\([a-\delta,a+\delta]\). It is false if issued and
\(|\theta-a|>\delta\). Phase 18 instead constructs a separate bound on
source-anchor error from nuisance bands. With estimated source tail rate
\(p_S\), target tail band [\(p_L\),\(p_U\)] and residual upper bound
\(q_U\), its rule requires
\(239.45\max(|p_L-p_S|,|p_U-p_S|)+239.55q_U\leq\delta\) and policy
admission. The target score interval uses the conditional mean bands and
the residual term separately. Returning full physical support on
rejection does not issue a certificate. Certificate issuance, false
certificates per run, false certificates among issued, and threshold
abstention have distinct denominators.

\begin{figure}[htbp]
\centering
\includegraphics[width=\linewidth]{figures/diagnostic.png}
\caption{Diagnostic framework. Top row follows the population interval and threshold call. Bottom row follows the anchor decision and two risk denominators. Boxes are conceptual, with no plotted observations.}
\label{fig:diagnostic}
\end{figure}

\section{Experimental design and comparators}

\begin{figure}[htbp]
\centering
\includegraphics[width=\linewidth]{figures/design.png}
\caption{Design and reproduction architecture. The upper row shows synthetic information and acquisition. The lower row shows the artifact and numerical verification path. This schematic has no measured quantities.}
\label{fig:design}
\end{figure}

Phase 18 supplies one common package of a fixed-allocation frame stream,
a probability-mixture audit stream, source anchor data and conditional
evaluation counts. Direct calibration uses the audit stream. Stratified
direct uses the frame. The combined rule uses the stratified band
subject to registry checks and a disagreement test. Pooling is available
only when declared design and common current-law assumptions justify it.
A refusal control returns full support. In shared blindness the
independent frame and audit streams sample the same stale law. The
evaluation law differs. Conditional body and tail means are estimated
using Hoeffding bounds and nuisance masses use Clopper--Pearson bands in
this design. The Phase 19 extension isolates known conditional score
means and exogenous omission restrictions.

Phase 25 uses nine frozen case laws and two budgets. Each paid policy
receives the same pilot of 128 independent draws per stratum. It
discards these draws from inference and acquires additional observations
from independent streams. Within each replicate, policies use common
prefixes of those inference streams to retain paired comparisons. A
replicate, rather than a stratum, draw or duplicated row, is the unit of
analysis. Duplicated blocks are billed by row and analyzed by
independent unit. Pilot and inference together exhaust the synthetic
tariff B. Frame enumeration, cap construction, and externally supplied
score means are outside this tariff. Operational field costs remain
unknown.

The width proxy minimizes \(\sum_h w_h\allowbreak/\sqrt{n_h}\).
Paid-pilot cost-adjusted Neyman minimizes \(\sum_h w_h^2
v_h\allowbreak/n_h\) using smoothed pilot variance \(v_h\). Minimax uses
the Bernoulli variance bound one-quarter. Proportional allocation
approaches \(w_h\) shares in independent units, while equal spend
divides charges evenly between strata. All use the same envelope. Direct
reachable calibration buys equal-spend inference and omits expansion. It
is a diagnostic under incomplete coverage. The true-rate variance oracle
has privileged information. A zero-acquisition abstention control is
reported separately. None of the latter controls defines a
matched-information superiority claim.

\section{Statistical analysis}

Phase 18 scores utility as +1 for a correct certificate, -1 for a false
one and 0 for none. Three paired combined-minus-comparator contrasts use
an equally weighted mixture of four restricted cases. Replicate
differences lie in \([-2,2]\). Two-sided Hoeffding intervals
\citep{hoeffding1963} use \(\alpha=.05/3\). The safety bound in the
shared-blind case is the one-sided Clopper--Pearson \citep{clopper1934}
upper bound for the unconditional false-certification proportion with
error .05/9. It belongs to the nine-case combined-policy family. It is
distinct from the conditional proportion among issued certificates. This
descriptive mixture supplies no deployment weights.

The Phase 25 primary comparison uses small omission at 32,768 charged
units. For independent paired replicates j let \(D_j\) be proxy
useful-precision indicator minus Neyman useful-precision indicator. Its
mean is the estimated difference. Since \(D_j\) lies in \([-1,1]\), the
radius is \(\sqrt{2\log(2/\alpha)/n}\) with \(\alpha=.05\), truncated at
the possible difference range. The separate secondary family uses
\(\alpha=.05/17\) per contrast. The primary improvement rule required
the lower endpoint to exceed 0.05. Equal spend comparisons are
descriptive here. Zero in an interval does not establish equivalence.
The exploratory sensitivity sweep uses a separate seed and one-factor
interventions, with no confirmatory ranking.

\section{Results}

\subsection{Shared blindness separates coverage and decision safety}

The combined policy covered 0/512 target means and issued 180
certificates. All 180 were false. The simultaneous upper unconditional
false-certification bound was 40.73\%. Wrong threshold labels numbered 0
because the missed intervals could remain on the correct side of the
threshold. The shared frame agreement was insufficient for the anchor
claim.

\begin{figure}[htbp]
\centering
\includegraphics[width=\linewidth]{figures/shared.png}
\caption{Phase 18 shared blindness. Percentages use the denominator named above each panel. The full-run false-certificate count and among-issued count share a numerator. They answer different risk questions.}
\label{fig:shared}
\end{figure}

\begin{table}[htbp]
\centering
\small
\begin{tabular}{lrr}
\hline
Endpoint & Count & Denominator \\
\hline
Mean coverage & 0 & 512 \\
Certificates issued & 180 & 512 \\
False certificates (all runs) & 180 & 512 \\
False certificates (among issued) & 180 & 180 \\
Wrong threshold labels & 0 & 512 \\
\hline
\end{tabular}
\caption{Frozen Phase 18 shared-blindness decisions. The one-sided simultaneous risk bound applies to the all-run false-certificate row.}
\label{tab:shared}
\end{table}

At equal spend in the separate Phase 19 large-omission case, the
supplied cap raised coverage from 0/512 to 512/512 on the same counts.
Useful precision was 0/512. The 1.171-unit identification floor exceeded
the precision limit.

\subsection{The original allocator lost the primary comparison}

At 32,768 charged units the width proxy gave 1/256 covered-and-precise
intervals, pilot Neyman 104/256 and equal spend 229/256. The primary
difference was -40.23 percentage points with 95\% interval [-57.21,
-23.26]. All 103 discordant pairs favored Neyman. All 17 adjusted
secondary intervals included zero. All three covered every target. The
per-replicate width cutoff therefore explains the useful-precision gap.
The improvement gate failed.

\begin{table}[htbp]
\centering
\small
\begin{tabular}{lrrp{0.19\linewidth}}
\hline
Allocation & Covered & Covered and precise & Mean width (score units) \\
\hline
Width proxy & 256/256 & 1/256 & 0.6726 \\
Pilot Neyman & 256/256 & 104/256 & 0.6113 \\
Equal spend & 256/256 & 229/256 & 0.5576 \\
\hline
\end{tabular}
\caption{Phase 25 primary case. Coverage and useful precision are counts per replicate. Mean width is a summary of complete score intervals, not a classification rule.}
\label{tab:allocation}
\end{table}

\begin{figure}[htbp]
\centering
\includegraphics[width=\linewidth]{figures/allocation.png}
\caption{Phase 25 primary comparison. Left shows total coverage and covered-and-precise percentages, with useful counts labeled. Right decomposes mean width in score units into weighted stratum bands and omission expansion. The dotted line is the per-replicate width limit, shown for scale.}
\label{fig:allocation}
\end{figure}

\section{Failure analysis and sensitivity}

In the primary case the proxy allocated fewer independent units to the
high-rate stratum. Its weighted contribution to the mean width was
0.4393 score units, versus 0.3702 for Neyman and 0.2953 for equal spend.
The corresponding envelope increments were 0.1168, 0.1169 and 0.1169.
For each replicate, the reachable width is \(234.2\sum_h w_h(u_h-l_h)\).
The envelope increment is the final score interval width minus that
reachable width. Averaging is linear for this decomposition, but the
event ``covered and precise'' is an indicator of two inequalities on a
single replicate. Mean components explain the observed width gap and
cannot reconstruct useful counts by comparing a mean with
\ensuremath{\tau}. The frozen audit checked integer allocation against
the continuous optimum and did not find rounding that could explain the
allocation gap.

Neyman's pilot responds to observed variance and shifted draws toward
the stratum driving width. Equal spend placed more draws there again
under the frozen prices, and led descriptively. Estimating variance from
a short pilot can be unstable when hits are rare. Smoothing prevents
zero variance estimates and can inflate very small rates. Even true
variance optimization does not optimize the tail-sensitive likelihood
width or useful precision. The comparison therefore establishes a local
objective mismatch under these costs and laws.

\begin{figure}[htbp]
\centering
\includegraphics[width=\linewidth]{figures/sensitivity.png}
\caption{Selected Phase 26 exploratory interventions. Cells show covered-and-precise counts over the number of synthetic runs. Selection spans pilot, price, omitted mass, cutoff and budget changes. These are descriptive.}
\label{fig:sensitivity}
\end{figure}

\begin{table}[htbp]
\centering
\small
\begin{tabular}{lrrr}
\hline
Exploratory setting & Width proxy & Pilot Neyman & Equal spend \\
\hline
Reference & 3/96 & 45/96 & 84/96 \\
Pilot 32 per stratum & 3/96 & 3/96 & 85/96 \\
Pilot 512 per stratum & 2/96 & 74/96 & 81/96 \\
Equal tariffs & 84/96 & 95/96 & 96/96 \\
Large omission & 0/96 & 0/96 & 0/96 \\
Width limit 0.50 & 0/96 & 0/96 & 3/96 \\
Width limit 0.75 & 96/96 & 96/96 & 96/96 \\
Budget 65,536 & 96/96 & 96/96 & 96/96 \\
\hline
\end{tabular}
\caption{Machine-readable extract of the same exploratory cells. The full frozen sensitivity grid also includes invalid weights, bound errors and prevalence changes.}
\label{tab:sensitivity}
\end{table}

Pilot size changes both variance information and remaining acquisition
budget. Large omission and strict width limits produced floor effects,
while higher budgets or a loose width limit produced ceilings. Invalid
weights and omitted-mass caps could destroy coverage across methods. An
invalid hidden-tail restriction could escape detection at small omission
because sampling bands remained wide. The separate large-omission
invalid-support case retained its negative result. These observations do
not provide a policy for diagnosing unknown field assumptions from
agreement alone.

\section{Reproducibility, validity and implications}

The public artifact pins source, configuration and archive members. The
separate Phase 27 implementation checks 36,864 Phase 25 method-replicate
records in 144 cells from counts, reconstructs endpoints and paired
intervals, and checks Phase 18 decisions from verified arrays and raw
source counts. The historical replay independently reconstructs the
Phase 18 arrays upstream. The reference implementation shares the
declared design, data and numerical environment. The project executed
the reproduction internally. The publication builder derives figures,
tables and interpolated result text from frozen machine-readable
evidence. The final audit also checks per-replicate width classification
and mean-width components from the Phase 25 records. Source and output
hashes identify the rendered PDF.

Construct validity depends on the chosen loss. Mean coverage, a correct
threshold decision, interval width, and certification each describe a
different property. External validity depends on knowing which frames
miss which units and whether the omitted component and costs resemble
these cases. There is no estimate of synthetic-case prevalence in a
field population. The design holds weights and score support fixed in
the primary comparison, assumes independence at the replicate level, and
uses externally supplied bounds rather than measured field bounds. Wider
confidence intervals, conservative refusal, and narrower credible
identification sets carry different operational consequences. The
benchmark can help distinguish these choices. It supplies no empirical
safety estimate or biological validation.

\subsection{Software and worked reproduction}
The \pkg{GENEVRA} research drivers are written in \proglang{Python}. The archival
runner verifies the capsule before replay. The independent reference reads
stored counts and independently verifies intervals and decisions. The publication
driver reads checked evidence and renders the manuscript assets. The software
guide documents the module interfaces, installation, and error handling.

From the repository root the following command rebuilds the Phase 29 review PDF
from frozen artifacts. It checks the output against the published bytes.
\begin{Code}
sh scripts/reproduce_final_research.sh --publication-only /tmp/genevra-review
\end{Code}
The software guide also gives the separate complete historical replay command.
Figures~\ref{fig:diagnostic} and~\ref{fig:design} describe the workflow.
Figures~\ref{fig:shared}, \ref{fig:allocation}, and~\ref{fig:sensitivity}
and Tables~\ref{tab:notation}, \ref{tab:shared}, \ref{tab:allocation},
and~\ref{tab:sensitivity} are produced from the existing publication inputs.

\section{Conclusion}

Shared sampling-frame omissions can defeat certification despite
agreement. Valid external information can restore coverage without
meeting a precision objective. In the matched small-omission comparison,
GENEVRA's width proxy underperformed both the variance-aware and
equal-spend allocations on the useful-precision endpoint. That negative
result is a test of the benchmark's ability to expose failures under
equal information and charges. The next evidence is independent artifact
review, a withheld comparison with locked policies for transfer claims,
and independently governed frame and outcome measurements for any
empirical claim. Journal submission remains on hold.

\bibliography{references}

\end{document}
